\documentclass{article}
\usepackage[utf8]{inputenc}
\usepackage[T1]{fontenc}
\usepackage{geometry}
\usepackage{graphicx}
\usepackage{float}
\usepackage{hyperref}
\usepackage[english,indonesian]{babel}

\title{Optimasi Pemrosesan Template LaTeX Menggunakan Pendekatan Berbasis AI}
\author{Academic Support Agent}
\date{\today}

\begin{document}
\maketitle

\begin{abstract}

\end{abstract}

\chapter{PENDAHULUAN}


\chapter{TINJAUAN PUSTAKA}


% Original content (commented out):
% %%% Modelo de Dissertação ou Tese do Programa de Pós-Graduação em Ciências Toxinologia (PPGTox) do Instituto Butantan.
%% Criado por Fernanda Midori Abukawa (fernanda.abukawa@gmail.com)

\documentclass[12pt]{article}

% Idioma Português
\usepackage[portuguese]{babel}

% Formatação da Página
\usepackage[a4paper,top=3cm,bottom=2cm,left=3cm,right=2cm]{geometry}

% Fonte Arial
\usepackage{helvet}
\renewcommand{\familydefault}{\sfdefault}
\newcommand{\bigsize}{\fontsize{16pt}{20pt}\selectfont}

% Espaçamento 1.5
\usepackage{setspace}
%\linespread{1.25}
\onehalfspacing

% Tamanha da fonte das seções 
\usepackage{sectsty}
\sectionfont{\fontsize{14}{15}\selectfont}
\subsectionfont{\fontsize{12}{15}\selectfont}
\subsubsectionfont{\fontsize{12}{15}\selectfont}
\paragraphfont{\fontsize{12}{15}\selectfont}


% Parágrafos
\usepackage{indentfirst}
\setlength{\parindent}{1.25cm}

% Numero das paginas
\usepackage{fancyhdr}
\pagestyle{fancy}
\fancyhead[R]{\thepage}
\fancyhead[L]{}
\fancyfoot{}
\renewcommand{\headrulewidth}{0pt}


% Pacotes
\usepackage{amsmath}
\usepackage{graphicx}
\usepackage[colorlinks=true, allcolors=black]{hyperref}
\usepackage{titlesec}
\usepackage{float}
\usepackage{pdfpages}
\usepackage{ragged2e}


% começar nova seção em nova página 
\newcommand{\sectionbreak}{\clearpage}

% ponto depois do número da seção 
\usepackage{titlesec}
\titlelabel{\thetitle.\enspace}

% Bibliografia
\usepackage[abnt-emphasize=bf]{abnt-alf}
\bibliographystyle{abnt-alf}

% Lista de Tabelas e Figuras 
\usepackage{tocloft}
\tocloftpagestyle{empty}
\renewcommand{\cftfigfont}{Figura }
%\renewcommand\cftloftitlefont{\hfill\large\bfseries\hfill} % optional
\renewcommand{\cfttoctitlefont}{\hspace*{\fill}\normalsize\bfseries}
\renewcommand{\cftaftertoctitle}{\hspace*{\fill}}
\renewcommand{\cftlottitlefont}{\hspace*{\fill}\normalsize\bfseries}
\renewcommand{\cftafterlottitle}{\hspace*{\fill}}
\renewcommand{\cftloftitlefont}{\hspace*{\fill}\normalsize\bfseries}
\renewcommand{\cftafterloftitle}{\hspace*{\fill}}
\renewcommand\cftfigaftersnum{:} % affix ")" after figure number
\renewcommand{\cfttabfont}{Tabela }
\renewcommand\cfttabaftersnum{:} % affix ")" after figure number

% renomear termos 
\renewcommand{\figurename}{Figura}
\renewcommand{\tablename}{Tabela}
%\renewcommand\refname{Referências}
\addto\captionsportuguese{\renewcommand{\refname}{REFERÊNCIAS$^1$}}
%\addto\captionsportuguese{\renewcommand{\bibname}{REFERÊNCIAS}}
\addto\captionsportuguese{\renewcommand{\contentsname}{SUMÁRIO}}
\addto\captionsportuguese{\renewcommand{\listfigurename}{\centering LISTA DE ILUSTRAÇÕES}}
\addto\captionsportuguese{\renewcommand{\listtablename}{\centering LISTA DE TABELAS}}

% Legendas 
\usepackage[labelsep=endash]{caption}
\captionsetup[figure]{justification=justified, singlelinecheck=off,labelfont=bf}
\captionsetup[table]{justification=justified, singlelinecheck=off,labelfont=bf}

%%%%%%%%%%%%%%%%%%%%%%%%%%%%%%%%%%%%%%%%%%%%%%%%%%%%%%%%%%%%%%%%%%%%%%%

\title{Your Paper}
\author{You}

\begin{document}
%\maketitle

\input{titlepage}


\newpage

\input{Resumo}

\newpage

%% Lista de ilustrações (opcional)
%\pagestyle{empty}
%\listoffigures

%\newpage

%% Lista de tabelas (opcional)
%\listoftables
%\clearpage
%\pagestyle{fancy}
%\newpage

%%% Sumário %%%%%%
\newpage
\pagestyle{empty}
\tableofcontents
\clearpage
\pagestyle{fancy}


%%%%%%% Introducao %%%%%%%%%%%%%%
\section{\normalsize{INTRODUÇÃO}}
\vspace{3mm}
\input{intro}

\newpage

%%%%%%% Obejtivos %%%%%%%%%%%%%%%%

\section{\normalsize{OBJETIVO}}

\subsection{Objetivo Geral}

Escrever objetivo geral.

\subsection{Objetivos Específicos
}
\begin{itemize}
    \item Objetivo específico.
\end{itemize}

\newpage

%%%%%%%% Mat&Met %%%%%%%%%%%%%%%%%%%%

\section{\normalsize{MATERIAIS E MÉTODOS}}
\vspace{3mm}
\input{mat&met}

\newpage

%%%%%%%% Resultados %%%%%%%%%%%%%%%%%%%%

\section{\normalsize{RESULTADOS}}
\vspace{3mm}

\input{resultados}

%%%%%% Discussão %%%%%%%%%%%%%%%%%%%%%

\section{\normalsize{DISCUSSÃO}} 

\input{discussao}

\newpage

%%%%%%% Conclusão %%%%%%%%%%%%%%%%%

\section{\normalsize{CONCLUSÃO}}

\input{conclusao}

\newpage

%%%%%%%%%%%%%%%%%%%%%%%%%%%%%%%%%%%%

\nocite{scikit-learn}
\titleformat*{\section}{\bfseries\normalsize\centering}
\addcontentsline{toc}{section}{REFERÊNCIAS}
\section*{}
\vspace{-1cm}
\footnotetext[1]{ASSOCIAÇÃO BRASILEIRA DE NORMAS TÉCNICAS. \textbf{NBR 6023}: Informação e documentação - referências - elaboração. Rio de Janeiro: ABNT, 2018.}

\begin{singlespace}

\bibliography{references}

\end{singlespace}

\newpage

\appendix
\renewcommand\thefigure{\thesection A\arabic{figure}} 
\section*{APÊNDICE A}\label{Apendice}
\addcontentsline{toc}{section}{APÊNDICES}
%\renewcommand\thefigure{\thesection.\arabic{figure}} 
\setcounter{figure}{0}

\input{apendice}

\newpage

%% para tabelas, opcional fazer um arquivo separado
%\setcounter{table}{0}
%\renewcommand{\thetable}{A\arabic{table}}
%\input{apendice_tabelas}

\end{document}

\end{document}